\documentclass[11pt,a4paper]{article}
\usepackage[margin=22mm]{geometry}
\usepackage{amsmath,amssymb,amsthm,hyperref}
\newtheorem{proposition}{Proposition}
\newtheorem{lemma}[proposition]{Lemma}
\title{Primitive maximal-parabolic reduction for the alternating bound}
\author{Human Compression project: L044 working continuation}
\date{}
\begin{document}\maketitle
This section reuses the accepted complete L047 primitive/common
theory. Its bounded cases are supplied by the complete required certificate
\texttt{primitive\_certificate.g}, rather than taken from the target
source's rank table.

\begin{proposition}
Let $G<A_n$ be primitive and have a string C-marking of rank $d$.
For every $n\geq12$, one has $d\leq\lfloor(n-3)/2\rfloor$.
\end{proposition}
\begin{proof}
If the noncommuting diagram is disconnected, the accepted commuting
factor lemma identifies $G=T\times T$, with $T$ nonabelian simple
and $n=|T|\geq60$. The intersection property gives
$2^d\leq |G|=n^2$, so $d\leq2\log_2n<(n-3)/2$.
The complete finite argument below settles $12\leq n\leq24$;
for the rest of this proof assume $n\geq25$.

For a connected diagram use the exact retained Mar\'oti alternative
as in L047, including $A_m^\ell\leq G\leq S_m\wr S_\ell$,
$m\geq5$, and $1\leq k\leq m/2$ after complementing subsets,
where $n=\binom mk^\ell$. If $\ell=1$, the case $k=1$ is excluded
by $G<A_n$. Whiston in the faithful $m$-point action gives
$d\leq m-1$. Since $n\geq12$, the remaining range has $m\geq6$,
and $2(m-1)\leq\binom m2-3\leq n-3$.
If $\ell\geq2$, the same abstract group has the faithful transitive
imprimitive action on $m\ell$ letters supplied in L047. Its accepted
bound gives $d\leq m\ell/2+1\leq(n-3)/2$, because
$m^\ell\geq m\ell+5$ for $m\geq5,\ell\geq2$.

The Mathieu actions all have degree at most24 and belong to the
finite boundary. In the remaining alternative put
\[
 M(n)=n\prod_{j=0}^{\lfloor\log_2n\rfloor-1}(n-2^j).
\]
If $d>\lfloor(n-3)/2\rfloor$, then $2d-1\geq n-3$ and $d\geq12$
for $n\geq25$. The retained Conder bound, used only in its valid
rank range $d\geq9$, gives $|G|\geq2^{2d-1}$.
Direct integer multiplication verifies
$M(n)<2^{2(\lfloor(n-3)/2\rfloor+1)-1}$ for $25\leq n\leq63$.
For $n\geq64$, write $q=\lfloor\log_2n\rfloor\geq6$. Then
\[
 M(n)<n^{q+1}<2^{(q+1)^2}\leq2^{2^q-3}\leq2^{n-3},
\]
where $(q+1)^2\leq2^q-3$ follows from the case $q=6$ by induction.
This contradicts the lower bound.
\end{proof}


\section*{The bounded primitive boundary}
For $12\leq n\leq24$ the required program runs through every
primitive catalogue action contained properly in $A_n$. The
retained completeness of GAP's primitive group library is explicit.
Its initial necessary bounds are these:
\begin{itemize}
\item an independent $d$-tuple gives a strict subgroup chain, so
$d\leq\Omega(|G|)$, with prime factors counted with multiplicity;
\item involutions must generate $G$;
\item an endpoint involution $t$ centralises the other $d-2$
nonneighbouring generators. By IP it is independent of them,
so $d\leq1+\Omega(|C_G(t)|)$;
\item the alternating-index generators form an elementary abelian
subgroup of rank $\lceil d/2\rceil$.
\end{itemize}
The actual degree12 $M_{11}$ action and degree13 $\operatorname{PSL}_3(3)$
action have elementary abelian two-rank two, verified by all subgroups
of a Sylow two-subgroup. This excludes ranks at least five. The
actual degree15 action is identified with $A_8$, so the retained
Whiston bound gives rank at most six. For degree24 $M_{24}$, the
involution centraliser orders are exactly $7680,21504$. At forbidden
rank $d\ge11$ its connected tail of rank $d-2\ge9$ satisfies
Conder's bound. The endpoint and tail intersect trivially by IP,
so its centraliser would have order at least
$2\cdot2^{2(d-2)-1}=2^{2d-4}\ge2^{18}$, a contradiction.

After these necessary tests the exact current GAP catalogue indices
are
\[
\begin{gathered}
 (12,2),\\
 (16,7),(16,10),(16,11),(16,12),(16,13),(16,15),\\
 (16,16),(16,17),(16,18),(16,19),(16,20).
\end{gathered}
\]
These are current library indices, not the source's Sims numbering.
The program verifies elementary abelian two-rank three for $M_{12}$,
so it need only exclude ranks five and six. Each listed degree16
socle is elementary abelian of order16. Its primitive affine action
has irreducible normal module $V=\mathbb F_2^4$, and the quotient
embeds in $\operatorname{GL}_4(2)\cong A_8$. The accepted L047
normal-module bound and Whiston theorem give
$m(G)\le m(G/V)+1\le7$. Thus rank seven is the only remaining
rank to exclude. No quotient intersection property is assumed.

\begin{lemma}[Completeness of the finite prefix search]
The required prefix search excludes every connected string
C-marking in the stated rank ranges. It does not use a
source-reported maximum rank.
\end{lemma}
\begin{proof}
List all involutions of the actual permutation group $G$. Starting
with the empty prefix, extend a prefix by every involution which
commutes with all earlier entries except its immediate predecessor
and does not commute with that predecessor. This is necessary for
a connected string diagram. Prefixes are considered up to the
normaliser of $G$ in the full symmetric group, then up to the
pointwise conjugation stabiliser of the already chosen prefix.
These are genuine symmetries of all completion conditions.

Every prefix of a C-marking is independent. It also satisfies
intersection of its initial and terminal maximal subgroups equal
to the middle subgroup; prefixes failing either test are rejected.
These are necessary tests, not an assertion that they alone imply
full IP. If a prefix subgroup $H$ has too small an index in $G$
to permit a strict chain of the remaining length, it is rejected.
There is one further necessary order test: if the target rank is
$t$ and the prefix has length $m$, the last $t-m-1$ future entries
centralise $H$ and generate a group intersecting $H$ trivially by
IP. They consequently embed in
$C_G(H)/Z(H)$, whose order must be at least $2^{t-m-1}$.

A prefix already generating $G$ cannot be extended independently.
On complete tuples the program tests all subset intersections
against their literal common-index subgroup. Thus every target
C-marking is reached, modulo the stated conjugations, unless a
necessary condition has already excluded it. The resulting complete
lists are empty for all twelve listed actions and rank ranges.
A positive control recovers an actual rank-three $A_5$ marking,
so the empty result is not an unconditional output of the engine.
\end{proof}

The complete prefix engine is inherited from the accepted L047
certificate, with only its name and returned report changed. Its
new rank ranges and the entire finite primitive screening are part
of the present required program. The thirty-nine integer Mar\'oti
inequalities for $25\le n\le63$ are verified there too.

\paragraph{Source scope and retained inputs.}
This proves the bound asserted in L044 Proposition~2.2 and hence
its Corollary~2.3: a primitive maximal parabolic in an $A_n$
C-marking has rank $r-1\le(n-3)/2$.
Mar\'oti's exact theorem (including its classification dependence),
Whiston, Conder in rank at least nine, the accepted L047 common
lemmas, and the primitive library and exact GAP group operations
remain declared premises. We do not reproduce the source's finer
individual maximum-rank claims in Table1 or attribute current
GAP catalogue indices to its historical Sims list.
\end{document}
